% ------------------------------------------------------------------
%  Prefazione (non numerata: \chapter nel frontmatter)
%  Scheda: ricerca/06_indice-ragionato.md, §6.1 (1.000–1.500 parole:
%    perché questo libro, da quale domanda nasce, perché comincia
%    dall'università, perché è gratuito).
%
%  STATO: prima stesura (08/10/2026), circa 1.300 parole.
%  - Nessun aneddoto nuovo: l'unico ricordo personale è quello già
%    scritto dall'autore nel cap. 1 (esame di ecofisiologia acquatica).
%  - Voce: «io» discreto; l'autore non si presenta come docente né come
%    psicologo (decisione 13).
%  - Nessun dato numerico: i dati sono nei capitoli.
%  - Le ragioni della scelta dell'università sono solo accennate: per
%    esteso sono nel cap. 26 («Dall'università alla scuola»).
%  - Licenza come nel colophon (CC BY-NC-SA 4.0); recapito come in A.8.
%  - Da ricontrollare a libro finito: data in calce; frase sulle norme
%    («autunno del 2026»); eventuale nota sull'uso di strumenti di IA
%    nella preparazione del libro (decisione dell'autore).
% ------------------------------------------------------------------
\chapter{Prefazione}

\iniziale{Q}{uesto libro} nasce da una domanda semplice: a che cosa
serve studiare, se poi si dimentica quasi tutto? Chiunque sia passato
per la scuola conosce la sensazione di aver superato una verifica o un
esame e, poche settimane dopo, di non ricordarne quasi nulla. Me la
sono posta anch'io, e non in astratto: nel primo capitolo racconto di
un esame preparato con cura, superato con il massimo dei voti e
dimenticato nel giro di pochi mesi. Per molto tempo ho pensato che
fosse normale, il prezzo inevitabile dello studio. In parte lo è:
dimenticare fa parte del modo in cui funziona la memoria. Ma
dimenticare quasi tutto, e così in fretta, non è un destino. È la
conseguenza di come si studia; e il modo di studiare si può cambiare.

La buona notizia è che il modo giusto è noto. Negli ultimi decenni la
ricerca sull'apprendimento ha chiarito, con centinaia di esperimenti in
laboratorio e in aula, che cosa fa ricordare a lungo e che cosa no:
richiamare dalla memoria invece di rileggere, distribuire lo studio nel
tempo invece di concentrarlo alla vigilia, alternare, spiegare,
collegare ciò che si studia a ciò che si sa già. Molte di queste cose i
maestri antichi le avevano intuite. Aristotele sapeva che gli esercizi
conservano la memoria \enquote{facendo ricordare di nuovo}; Quintiliano
si stupiva di quanta solidità aggiunga a ciò che si è imparato una
notte di sonno. La scienza di oggi ha messo alla prova queste
intuizioni, ha scartato ciò che non reggeva e ha aggiunto ciò che gli
antichi non potevano sapere.

Il problema è che questo sapere arriva poco a chi studia. Resta nelle
riviste scientifiche, quasi tutte in inglese. Nelle scuole e nelle
università italiane, salvo eccezioni, nessuno insegna come si studia:
lo si dà per scontato, come se fosse una cosa che si impara da sé. E
negli scaffali abbondano i metodi che promettono di imparare tutto in
pochi giorni. Così studenti seri e volenterosi passano ore e ore sulle
strategie che rendono meno --- rileggere, sottolineare, ripetere tutto
la sera prima --- proprio perché sono quelle che danno la sensazione
più convincente di sapere. Questo libro prova a colmare quella
distanza: a raccogliere in italiano ciò che la ricerca ha chiarito, con
le sue fonti e senza trucchi, e a tradurlo in pratiche quotidiane alla
portata di chiunque.

Non scrivo da psicologo né da insegnante. Sono uno che ha studiato, a
scuola e poi all'università, e che a un certo punto ha voluto capire
perché ricordasse così poco di ciò che aveva studiato bene. Per trovare
una risposta ho cercato di fare ciò che il libro, pagina dopo pagina,
ti chiederà di fare: non fidarmi della sensazione di sapere, e andare a
vedere. Ho cercato gli studi originali, e non soltanto i loro
riassunti; ho controllato i dati e le citazioni sulle fonti, e dove una
fonte era più debole o solo indiretta l'ho detto. Ho letto gli antichi
nella loro lingua, e le traduzioni dei passi citati sono mie, salvo
diversa indicazione. Soprattutto, ho cercato di distinguere sempre, con
parole semplici, ciò che è solido, ciò che è soltanto promettente e ciò
che è falso, anche quando è molto diffuso.

Il libro comincia dall'università. È lì che si studia più da soli, con
più materiale e meno guida; è lì che è stata condotta gran parte degli
esperimenti di cui parlerò; ed è lì che si formano gli insegnanti di
domani, attraverso i quali ciò che qui si dice può arrivare alla
scuola. Le ragioni di questa scelta, e il libro che dovrà seguire ---
un'edizione per le medie e le superiori, con i loro insegnanti e le
loro famiglie --- sono spiegate nel capitolo~\ref{cap:formatori}. Chi
è ancora alle superiori, e soprattutto chi si prepara all'esame di
maturità, troverà comunque qui il metodo intero, e un capitolo dedicato
al passaggio dalla scuola all'università
(capitolo~\ref{cap:transizione}).

I lettori a cui mi rivolgo sono due. Il primo è chi studia: dalla
matricola al laureando, e chi continua a studiare dopo la laurea, per
un dottorato, un concorso o il proprio lavoro. Il secondo è chi
insegna all'università, e chi forma gli insegnanti o si prepara a
diventarlo. Al primo parlo per quasi tutto il libro; al secondo è
dedicata l'ultima parte, che si può leggere anche da sola. Ma chi
insegna farebbe bene a non saltare le altre: non si può insegnare un
metodo di studio che non si conosce, e quasi nessuno lo ha insegnato,
a suo tempo, nemmeno a chi oggi sale in cattedra. Con entrambi uso il
\enquote{tu}. È il modo più diretto di parlare a chi legge e, per un
libro che chiede di fare e non soltanto di leggere, anche il più
onesto.

Questo libro è gratuito. Il PDF si scarica senza pagare nulla da
\url{https://www.lucalevi.com}; l'edizione a stampa ha un prezzo
contenuto e non è pensata per guadagnare. È un libro di impegno
sociale, nato dalla convinzione che saper studiare non debba dipendere
dalla scuola che si è avuta la fortuna di frequentare, né dai soldi per
le ripetizioni o per un corso di metodo. Per la stessa ragione è
distribuito con una licenza Creative Commons (BY-NC-SA 4.0): puoi
copiarlo, stamparlo, fotocopiare le schede dell'Appendice~\ref{app:schede}
per i tuoi compagni o per i tuoi studenti, adattarlo e tradurlo,
purché tu ne citi la fonte, non lo faccia per guadagno e rilasci con la
stessa licenza ciò che ne ricavi.

Conviene anche dire che cosa questo libro non è. Non è un manuale di
trucchi: il metodo che propone funziona, ma chiede fatica, e il
capitolo~\ref{cap:dueforze} spiega perché quella fatica è un buon
segno. Non è un testo clinico né un parere legale: il capitolo sugli
studenti con DSA e altri bisogni speciali
(capitolo~\ref{cap:dsa}) non sostituisce una diagnosi né il servizio
del tuo ateneo. E non è un libro definitivo. La scienza
dell'apprendimento continua a lavorare; le norme sull'università e
sulla formazione degli insegnanti cambiano spesso, e quelle citate sono
le norme in vigore mentre scrivevo, nell'autunno del 2026; le prove
sull'intelligenza artificiale (capitolo~\ref{cap:ia}) invecchiano nel
giro di pochi mesi. Ho cercato di non dire mai più di quanto le prove
permettano. Dove ho sbagliato, o dove qualcosa è cambiato, scrivimi:
\texttt{luca@lucalevi.com}. Le correzioni entreranno nel PDF e nelle
edizioni successive.

Un'ultima cosa. Questo libro non si limita a parlare del metodo: prova
ad applicarlo a se stesso, e ti chiederà di fare lo stesso. Ti chiederà
di chiuderlo spesso, di rispondere senza guardare, di tornare su
capitoli che credevi di sapere. Come, e perché, lo spiegano le pagine
che seguono. È più faticoso che leggere e basta; ma ciò che si impara
così non si perde dopo l'ultima pagina. Ed è proprio questo, in fondo,
che vuol dire studiare per la vita.

\begin{flushright}
  \itshape Cormons, ottobre 2026
\end{flushright}

\finecapitolo
